\section{Discussion}

To our knowledge, this paper presents the first exploration of BLV people’s experience with an AI guide in VR. This work successfully addresses our two overarching research questions by providing key insights into both the utilization (RQ1) and behaviors (RQ2) of BLV users towards an AI guide. Regarding RQ1, participants successfully used the guide to explore and gather reliable information about complex virtual environments. Some were even able to lead sighted groups on tours within those environments, effectively balancing navigation with social interaction. Regarding RQ2, we found that participants altered their treatment of the guide according to social contexts. While alone, they treated the guide as a tool to address their visual needs, keeping interactions brief and direct. However, in social situations, participants became more friendly and role-played with the guide according to its persona.

It should be noted that the behaviors we observed likely reflect a complex interplay between the guide’s design, social presence, and technical or environmental factors (e.g., system latency encouraging name-calling or “attention-grabbing” behaviors from users). Such behaviors offer a compelling foundation for future controlled experiments to hypothesize about and isolate the specific variables behind users’ actions towards an AI guide. Ultimately, these observations and our other insights are critical for developing future AI guides that are not only functionally effective but socially adaptable, enhancing accessibility and user experience for BLV individuals in diverse virtual environments.

\subsection{Comparing Human and AI Guides} \label{sec: Comparing Human and AI Guides}

We compared our results to our previous work on human guides in VR \cite{collins2023guide}, noting interesting differences between experiences with human and AI guides. While participants in both studies formed emotional connections with their guides, our AI-guided participants also displayed novel, AI-specific uses, suggesting a need for new features exclusive to AI guides.

We discuss our comparison with some caveats. For example, our human guide study did not include social situations where participants used the guide around confederates. Further, while people are typically experienced with human conversation, they are far less experienced interacting with an AI. Lastly, although participants in both studies were first-time guide users, using the human guide was closer to real-world conversation, whereas the AI guide required additional training for features like a button interface and audio beacons. These factors may have affected how users adapted to the guides and employed them in the studies.

Regardless of differences, both studies demonstrate that forming emotional connections with guides can be an enjoyable part of guidance. Users of the human guide often invited their guides to perform activities in VR with them, while users of the AI guide role-played with the guide to incorporate it playfully into their experience. These behaviors reflect prior research on sighted guides \cite{guerreiro2019cabot, ball2022runners}, assistive robots \cite{bonani2018robots}, and existing work that has shown people tend to form relationships with AI \cite{zimmerman2023humanai, gur2025algorithm}. Participants leveraged features like the inviting avatars of the virtual guides to develop a degree of closeness, though a lack of other desired features—like the guide dog reacting to being petted (Sara)—limited this connection.

We also observed differences in how participants used the AI guide compared to a human one. Beyond its primary function as a navigational assistant, the AI guide was used as a \textbf{general knowledge assistant}. Participants frequently asked simple, factual questions like, "What is a gazebo?" This contrasts with the human guide, which was used as a \textbf{technical support specialist} and asked more complex questions about the VR system and technical issues. These different uses are revealing: participants saw the AI as a reliable source for simple information, a digital tool for rote inquiries that they would not want to burden human guides with.

Finally, we observed key differences in how participants viewed the roles of human and AI guides during error prevention and recovery. Participants in the human guide study expected their guide to be \textbf{proactive}, offering preemptive warnings to avoid virtual hazards. In contrast, participants treated our AI guide as a \textbf{reactive failsafe}. They did not ask it to prevent mistakes but instead used it to recover from them, such as by asking it to teleport them to a safe location after becoming lost. This suggests that users hold AI guides to a different standard, expecting them to be tools for recovery rather than prevention.

Based on our findings, we propose the following design recommendations for virtual guides:

\begin{itemize}
\item \textbf{Foster Emotional Connections}: Future AI guides could incorporate interactive capabilities (see section \ref{sec: Advantages of Embodied, Personable Accessibility Assistants} to respond appropriately to role-playing.
\item \textbf{Encourage Advanced Usage}: Future AI guides should incorporate features that explicitly demonstrate their capacity for complex tasks. This could include having the guide proactively offer visual details or check on users who appear to be struggling.
\item \textbf{Align with User Roles}: Future AI guides should offer explicit failsafe mechanisms. Rather than having users give voice commands requesting the guide to teleport them to certain locations when they get stuck or lost, users could perform a simple button press that instantly returns them to a safe, known location.
\end {itemize}

\subsection{Advantages of Embodied, Personable Accessibility Assistants} \label{sec: Advantages of Embodied, Personable Accessibility Assistants}

Prior approaches to VR accessibility have focused on informative tools activated through precise controls, similar to screen readers. For example, Microsoft’s Scene Weaver \cite{balasubramanian2023sceneweaver} conveys information about avatars and interactable objects through customizable haptic and audio cues, or manually-triggered alt text. In contrast, many commercial VR applications employ virtual assistants imbued with personality \cite{nathie2019presence, oculus2019first, moss2018review, ven2021vr, cosmonious2022walkthrough}. These assistants are often comical or endearing companions guiding users through tutorials, enhancing engagement and fostering emotional connections between user and assistant \cite{nathie2019presence, cosmonious2022walkthrough, lakereddit, fink2022review, graf2022emotional, clavel2013artificial, sutskova2023cognitive, freeman2025comforting}.

Despite their success in mainstream VR, such personality-infused assistants have not been widely applied in accessibility contexts. However, emerging research indicates potential benefits \cite{phutane2023speaking}. Phutane et al. explored personality-driven conversational assistants for BLV users, reporting positive responses to their interactivity \cite{phutane2023speaking}. Building on this, our AI guide is an embodied agent with personality, assisting users through conversation rather than complex controls. Our observations of participant-guide interactions revealed several advantages of this approach. Namely, the guide leveraged users’ natural desire to role-play with assistive technology, which allowed them to rationalize AI mistakes, reduce embarrassment, and encourage positive social interactions. For example, when the guide appeared in its dog form, confederates pet it, mirroring the social benefits reported by guide dog handlers in the physical world \cite{hwang2024towards}.

Future AI guide designs should support role-playing behaviors beyond mere speech modifications to enhance the user experience. For instance, an AI guide in dog form could play fetch, while a robot form could interact with other technologies. Future studies should also explore BLV users' reactions to various guide embodiments, as we found the guide's appearance significantly influenced the extent of user role-playing. Future research could explore how users respond to guides resembling characters from popular media or traditional assistive devices, like white canes. Such research could also improve how users learn to work with guides. For instance, using a well-liked dog form in tutorials may foster patience during the learning process.

\subsection{Improving Interactions with AI Guides} \label{sec: Improving Interactions with AI Guides}

Although participants had been told they could frame their questions to the guide in any way they wanted, participants tended to re-use simplistic orders (e.g., “Take me to X”, “Describe what’s in front of me”). For example, although several participants disliked the guide’s inconsistent manner of describing their surroundings, none of them tried prompting the guide to give descriptions in specific ways.

This reflects current work on how non-expert users work with modern-day LLMs \cite{zamfirescu2023johnny, brachman2024knowledge, ha2024clochat}. Non-experts typically over-generalize their instructions, and do not provide explicit guidance on how they want the LLM to respond \cite{zamfirescu2023johnny}, leading to less-than-ideal output. 

However, the design of the AI guide's personas may have misled participants as well. For example, when Len told the guide, "Let's go smell the flowers," he likely expected the talkative guide to remember the flower patch it had just mentioned. The guide's verbosity may have implied a conversational understanding and memory that it did not possess. To better align user expectations with the guide's abilities, future versions could include having the guide's appearance reflect its capabilities, such as a small puppy guiding the user rather than a full-grown dog, to indicate inexperience and encourage users to “train” it with more specific commands. The guide could also encourage users to be more explicit with their queries, proactively teaching them how to use it most effectively.

\subsection{Limitations and Future Work} \label{sec: Limitations and Future Work}

Our study was conducted during a period of rapid evolution in generative AI. We started our prototyping and user studies in mid-2023, shortly after the release of GPT-4. While the landscape of AI tools has changed drastically since then—with new APIs offering superior performance and the public’s familiarity with LLMs having grown exponentially—our work demonstrates that prototyping with contemporary tools is necessary to generate lasting insights into human-AI interaction.

Despite its value, our guide prototype faced several architectural limitations due to the systems available at the time. One such drawback was the guide's delayed response time. This delay is an inherent challenge when utilizing external, cloud-based LLMs like GPT-4, as the processing pipeline involves (1) capturing and transmitting user audio and environmental snapshots, (2) processing the query with GPT-4, and (3) generating and streaming the audio response back to the VR headset. Our delays were further exacerbated because, at the time of our study, we lacked access to real-time response streaming. We were forced to wait for the entire response to be generated before transmitting the complete audio file, a constraint that would be largely mitigated by modern streamable Text-to-Speech APIs (e.g., PlayHT \cite{playhtapi}).

Another architectural limitation was the guide's low accuracy (63.2\% across all queries). We attributed many inaccurate responses to technical issues, including the misinterpretation of participant accents, incomplete user queries, or unresponsive APIs. Modern LLMs and Speech-to-Text models offer greater robustness to accent differences. Furthermore, modern APIs often allow for including the context of past queries in LLM calls, which would significantly improve the guide's ability to infer appropriate responses to vague or incomplete prompts. Introducing these features into a modern AI guide architecture would likely enhance its accuracy. Future work should examine which types of guide failures still occur despite these modern advancements and how they impact the user experience, as seen in prior work \cite{goetsu2020different}.

Beyond technical limitations, our study found that participants primarily used a naive approach to prompting the guide. At the time of our study, the general public was becoming aware of AI tools like ChatGPT, but usage was still not widespread. Users today are likely more experienced with prompt engineering and adapting their language to an AI’s perceived capabilities. Future work should include a longitudinal study to examine how users build expertise in crafting sophisticated queries, such as "Describe what’s in front of me using cardinal directions," and how this expertise impacts their perception of the guide's effectiveness. To support this evolving user base, we propose the following recommendations for future AI guides that can help new and experienced users craft effective queries:

\begin{itemize}
\item \textbf{Encourage Advanced Usage}: Guides should actively assist users in prompting them, offering examples of better prompts or teaching users prompt-engineering best practices to help users move beyond basic queries faster.
\item \textbf{Align Personas with Capabilities}: The guide's persona should visually and verbally reflect its capabilities. For instance, a robot persona might encourage blunt, direct language, while a human-like guide could inspire complex commands from users seeking a conversational expert.
\end{itemize}

As AI continues to rapidly evolve and render older implementations obsolete, it remains crucial and valuable to prototype with these developing tools. Findings generated about user behavior and interaction patterns with AI often possess a longevity that the underlying technology lacks. Research on human-AI interaction, including studies on simple voice assistants like Amazon’s Alexa, established enduring trends like the tendency to anthropomorphize AI or rely on it for emotional support \cite{lopatovska2018, gao2018}. Similarly, we expect the user behavior findings from our study—such as the tendency to roleplay with AI assistants or alter communication styles in solitary versus social settings—to carry forward. These findings can inform the design of future, more powerful AI tools, regardless of the architectural advantages that older prototypes might miss.